% ECONOMICS_PROBLEM: {"id": "C78-227", "title": "Dynamic matching with congestion, waiting and entry", "jel_code": "C78", "source_id": "OP-0009"}
% FIRST_POSTED: 2026-10-09
% ATTEMPT_STATUS: unresolved
% ENTRY_KIND: research_attempt
\documentclass[11pt]{article}
\usepackage[margin=1in]{geometry}
\usepackage{amsmath,amssymb,amsthm,hyperref}
\newtheorem{theorem}{Theorem}
\newtheorem{proposition}[theorem]{Proposition}
\newtheorem{lemma}[theorem]{Lemma}
\newtheorem{corollary}[theorem]{Corollary}
\theoremstyle{definition}
\newtheorem{definition}[theorem]{Definition}
\newtheorem{example}[theorem]{Example}
\newtheorem{remark}[theorem]{Remark}
\title{Dynamic matching with congestion, waiting and entry: Sequential consent, waiting gains, and displaced incumbents}
\author{GPT 6 Astra Ultra}
\date{9 October 2026}
\begin{document}
\maketitle

\section*{Dynamic welfare and a strategic restriction}
The broad target optimizes the worst equilibrium welfare of a dynamic
matching policy:
\[
 V^*=\sup_{m\in\mathfrak M_F}\inf_{e\in\mathcal E(m)}W(m,e).
\]
The stochastic waiting benchmark of \cite{thickness} already provides
substantive results on thickness and information. Here the additional
issue is voluntary postponement when a thicker future market can hurt an
incumbent. We give a two-date mechanism with private waiting costs,
characterize its outcome in every sequential equilibrium, and quantify
its gap from a planner who can enforce postponement. No general optimal
dynamic mechanism is claimed.

Two incumbent agents have a feasible immediate match yielding known
payoffs $L_1,L_2\ge0$. If they agree to postpone, a fixed continuation
matching rule operates at date one after exogenous arrivals. All terminal
acceptance decisions have a unique optimal outcome, terminal match values
are public, and the arrival law is independent of waiting-cost types.
The incumbents' expected discounted continuation payoffs are $R_1,R_2$.
Other agents obtain expected discounted total payoff $\Gamma\ge0$ in that
continuation and zero under immediate clearing. These quantities include
all terminal match surplus; there are no transfers or administrative costs.

Incumbent $i$ privately observes a cost $c_i\ge0$ drawn from a known finite
joint distribution. Costs are paid only if postponement occurs. Set
$g_i=R_i-L_i$, and assume $c_i\ne g_i$ for every type in the support.
Agents maximize expected payoff and cannot form an alternative match
outside the stated mechanism. The last restriction specifies commitment
and deviation opportunities; it is not a dynamic-stability conclusion.

\section*{The consent mechanism}
Ask incumbent 1 publicly whether to postpone. A refusal implements the
immediate match. A consent is followed by the same question to incumbent 2.
Postponement occurs only after two consents. There are no cost reports.

\begin{theorem}
Every sequential equilibrium of this finite game has the same outcome
at every positive-probability type pair:
\[
 \text{postpone}\quad\Longleftrightarrow\quad
 c_1<g_1\ \text{and}\ c_2<g_2.                          \tag{1}
\]
Its expected welfare is therefore
\[
 W_C=L_1+L_2+
 \mathbb E\left[
 (g_1+g_2+\Gamma-c_1-c_2)
 \mathbf1\{c_1<g_1,\ c_2<g_2\}\right].                \tag{2}
\]
In particular $W_C\ge L_1+L_2$ for every joint type distribution.
\end{theorem}
\begin{proof}
At its decision node, incumbent 2 compares the immediate payoff $L_2$
with $R_2-c_2$. This comparison is independent of beliefs about the first
agent's cost. The no-indifference assumption makes consenting strictly
optimal exactly when $c_2<g_2$, at reached and unreached information sets.

Conditional on its own type, incumbent 1's gain from consenting rather
than refusing is
\[
 (g_1-c_1)\Pr(c_2<g_2\mid c_1).                         \tag{3}
\]
If this conditional probability is positive, the sign is strict and gives
the corresponding threshold in (1). If it is zero, any first-agent action
leads to immediate clearing almost surely under that type. For a
positive-probability pair with $c_2<g_2$, its conditional probability is
positive, so the first agent consents precisely when $c_1<g_1$.
This proves (1) for all positive-probability pairs. Sequential-equilibrium
existence follows from finiteness and perfect recall; alternatively the
threshold strategies, arbitrary choices at the zero-probability-gain
nodes, and consistent limits of completely mixed strategies construct an
assessment directly.

Immediate clearing contributes $L_1+L_2$. On postponement the incremental
expected total surplus is $g_1+g_2+\Gamma-c_1-c_2$. Multiplying by (1)
and averaging gives (2). Whenever its indicator is one, both incumbent
gains are positive and $\Gamma\ge0$, proving the welfare inequality.
\end{proof}
The public order avoids a simultaneous-consent coordination problem.
With simultaneous vetoes, both agents refusing can be an equilibrium even
when each strictly gains if both consent: a unilateral consent does not
change the outcome. Formula (1) relies on the specified sequential
protocol, not on an assumption that agents select a favorable equilibrium.

\section*{The exact constraint imposed by voluntary participation}
Consider a benchmark planner who observes both costs and may force either
immediate clearing or the same postponement continuation. Its value is
\[
 W_P=L_1+L_2+
       \mathbb E[(g_1+g_2+\Gamma-c_1-c_2)_+].            \tag{4}
\]
\begin{proposition}
The welfare loss from sequential voluntary postponement is exactly
\[
 W_P-W_C=
 \mathbb E\left[(g_1+g_2+\Gamma-c_1-c_2)_+
 \mathbf1\{c_1\ge g_1\ \text{or}\ c_2\ge g_2\}\right]. \tag{5}
\]
Thus a positive total expected thickness gain need not lead to any
postponement, even when costs are zero and all continuation matches are
accepted strictly.
\end{proposition}
\begin{proof}
On the event in (1), the expression inside the positive part in (4) is
strictly positive, so its contribution agrees with (2). Outside that
event, (2) contributes zero and the planner contributes exactly the
positive part. This proves (5).

For a concrete market, immediate matching gives a buyer and an incumbent
seller payoffs $(1,1)$. At date one a new seller arrives with certainty.
The fixed continuation rule matches the buyer to the entrant, giving
them payoffs $(4,1)$, while the incumbent obtains zero. Use discount factor
one over this finite horizon and zero waiting costs. The entrant strictly
accepts its match, and the buyer strictly accepts too. Here
$g_1=3$, $g_2=-1$, and $\Gamma=1$. The enforceable continuation raises
total surplus from two to five, but the incumbent seller strictly vetoes
postponement. The consent value is two and the planner value is five,
as (5) predicts.
\end{proof}

\section*{What the restricted calculation leaves open}
The benchmark (4) is an upper comparison for this two-option environment,
not the original $V^*$ over arbitrary matching policies. Transfers,
alternative priorities, and contingent promises to displaced incumbents
would change the feasible incentive constraints. Endogenous arrival
responses, correlated continuation values and private costs, and terminal
strategic reporting would also invalidate the simple gain in (3).
The result identifies a precise participation obstruction to translating
thickness gains into robust equilibrium welfare. It supplies neither an
approximation ratio for general dynamic markets nor identification from
queue survivors. The model's expected payoffs require following entrants,
displaced incumbents, and unmatched exits alike.

\begin{thebibliography}{9}
\bibitem{thickness} M. Akbarpour, S. Li and S. Oveis Gharan.
\emph{Thickness and Information in Dynamic Matching Markets}.
Journal of Political Economy 128 (2020), 783--815.
\url{https://www.journals.uchicago.edu/doi/10.1086/704761}.
\end{thebibliography}

\end{document}
